% ===========================================================================
%  sviluppatore/01-introduzione/04-come-leggere.tex
% ===========================================================================
\section{Da dove cominciare, secondo cosa si deve fare}\label{sez:percorsi}

Questa guida non va letta tutta. Quattro percorsi, a seconda del motivo per
cui si è qui.

\begin{center}
\begin{tabular}{@{}L{48mm}L{98mm}@{}}
\toprule
\textbf{Si deve\ldots} & \textbf{Leggere} \\
\midrule
\rigaTenue
far partire una copia e
guardarsi intorno
  & \vediCap{ambiente}. Un'ora, e si ha il sistema che gira. \\

aggiungere una strategia,
un overlay, una pesatura
  & \vediCap{motore} per capire che cosa decide che cosa, poi il
    \vediCap{estendere}, che è una procedura da seguire con il codice
    aperto. \\

\rigaTenue
toccare l'interfaccia
  & \vediCap{dashboard}. Prima però \vediSez{strangler}: c'è un motivo
    se ci sono due applicazioni web e non una. \\

mettere mano a pagamenti,
abbonamenti, marketplace
  & \vediCap{marketplace}, e la parte sui permessi nel
    \vediCap{architettura}. È l'area dove un errore costa denaro vero. \\
\bottomrule
\end{tabular}
\end{center}

\begin{pratica}{Il minimo indispensabile}
Chi ha fretta e deve comunque toccare il codice legga due sole cose: la
separazione fra decisione ed esecuzione (\vediSez{decisione-esecuzione}) e le
due forme del sistema a plugin (\vediSez{due-forme}). Sono i due punti in cui
un errore non dà errore: produce un numero sbagliato, o un artefatto che non
compare.
\end{pratica}

\section{Che cosa questa guida non copre}\label{sez:fuori-perimetro}

\begin{itemize}
  \item \textbf{Il contenuto finanziario delle strategie.} Perché il momentum
        12-1 funzioni, se il PEAD sia ancora sfruttabile, come si tari un
        \emph{vol target}: sono domande di merito, e questo è un documento di
        ingegneria. Dove una scelta numerica è codificata nel sistema, la
        guida dice dov'è e che cosa cambia, non se sia giusta.
  \item \textbf{L'uso dell'applicazione.} Dove si clicca per lanciare un
        backtest sta nel \emph{Manuale utente}, documento separato nella
        stessa cartella.
  \item \textbf{Il codice del sito, riga per riga.} Il sito è documentato qui
        per ciò che espone alla metà locale: endpoint, formato dei token,
        regole dei permessi. Il resto è un'applicazione Next.js ordinaria.
\end{itemize}
